\chapter{Chapter 1 --- Introduction}\label{chapter-1-introduction}

\section{1.1 Context}\label{context}

Defects are expensive in proportion to how long they survive. A fault
caught at review costs a developer's attention; the same fault caught in
production costs an incident, a rollback and the trust of whoever
depended on the system. This asymmetry motivates \emph{defect
prediction}: statistical models that direct limited inspection effort
toward the code most likely to be wrong.

Classical defect prediction operates at release granularity, predicting
which files or modules will contain defects. Its practical difficulty is
one of timing. A prediction about a file, delivered at release, arrives
long after the change that introduced the problem, and often reaches a
developer who did not write it.

\textbf{Just-in-time software defect prediction (JIT-SDP)} moves the
unit of prediction from the file to the \emph{commit}. Each change is
scored as it arrives, so a risky change can be flagged while its author
still has it in mind. Kamei et al. \citep{kamei2013jit} established the
modern form of the task, defining a set of change-level metrics ---
size, diffusion, history, and developer experience --- that remain the
standard feature set, and reporting roughly 68\% accuracy and 64\%
recall at predicting defect-inducing changes.

The setting is unusual in three ways that jointly shape everything in
this report. Commits arrive as a \textbf{stream}, so any realistic
evaluation is temporal. Defect-introducing commits are \textbf{rare} ---
8.54\% on the corpus studied here --- so the task is severely
imbalanced. And, critically, \textbf{the label for a commit is not
available when the commit arrives.} It becomes available only when
someone finds the defect, traces it to its origin, and fixes it.

\section{1.2 The ground-truth problem}\label{the-ground-truth-problem}

JIT-SDP requires a label for every commit: did this change introduce a
defect? Almost no project records this directly. What projects record is
which commits \emph{fixed} defects.

The gap is bridged by the \textbf{SZZ algorithm} (\'{S}liwerski, Zimmermann
and Zeller), \citep{sliwerski2005szz}, which works backwards: take a
bug-fixing commit, identify the lines it modified, and use
\texttt{git\ blame} to find which earlier commit last touched each of
those lines. Those earlier commits are labelled defect-introducing.

The inference is only as good as its premise --- that the lines a fix
modifies are the lines that were wrong. \textbf{That premise fails
routinely, because bug-fixing commits are tangled.} Developers reformat
while fixing, rename while fixing, refactor while fixing, update
comments and tests while fixing. Every such line is blamed by SZZ, and
every commit that last touched one is labelled as having introduced a
defect it had nothing to do with.

The magnitude is now measured rather than suspected. Herbold et al.
\citep{herbold2022tangling}, using four annotators per line, found that
\textbf{between 17\% and 32\% of all changes within bug-fixing commits
actually address the underlying problem} --- rising to 66--87\% when
only production code files are considered. The remainder is tangled.
Roughly 11\% of lines were hard enough to classify that annotators
actively disagreed.

A substantial body of work has proposed refinements --- AG-SZZ, MA-SZZ,
R-SZZ, L-SZZ, RA-SZZ --- each adding filters intended to remove the
over-collection. Whether those refinements improve the labels, and at
what cost, is the subject of Chapter 4.

\textbf{The consequence for the field is structural.} Essentially every
JIT-SDP dataset in use descends from SZZ. Models are trained on SZZ
labels, evaluated against SZZ labels, and compared with one another on
SZZ labels. If those labels are systematically wrong, the entire
measurement apparatus is calibrated against a biased instrument --- and,
because the evaluation uses the same instrument as the training, the
bias is invisible from inside.

\section{1.3 The realism problem}\label{the-realism-problem}

A second problem is independent of labels and compounds them.

The standard evaluation protocol in the JIT-SDP literature is random
\emph{k}-fold cross-validation: shuffle the commits, hold out a fold,
train on the rest. On temporally ordered data this permits a model to
train on 2020 commits and be tested on 2015 ones. The resulting estimate
answers a question nobody has: how well would this model do if it could
see the future?

\textbf{Verification latency} is the subtler half. Even under a correct
chronological split, the assumption that a commit's label is available
at training time is false. The label exists only once the defect is
found and fixed, and on this corpus the median delay is 113 days, with
\textbf{53\% of defect labels arriving after a 90-day decision window.}

This is worse than a delay. A system with a 90-day window does not
simply wait --- it concludes that an unfixed commit is clean, trains on
that conclusion, and is corrected months later. \textbf{A learner under
verification latency is actively trained on labels it will subsequently
discover were wrong}, and the wrong labels it receives are exactly the
defect-introducing commits it most needs to learn from.

Cabral and Minku \citep{cabral2019orb} characterised this setting,
showing that the class imbalance itself evolves over time under latency,
and proposed \textbf{Oversampling Rate Boosting (ORB)} --- an online
ensemble that adjusts its resampling rate when predictions become biased
toward one class. ORB is the online learner used throughout this report.
The mechanism it uses to handle imbalance --- amplifying minority-class
examples in proportion to how rare they appear --- is also, on the
evidence assembled here, a plausible route by which label noise could do
disproportionate damage. Testing that is the next phase of the work and
is not attempted here.

\section{1.4 Problem statement}\label{problem-statement}

\begin{quote}
Just-in-time defect prediction is evaluated almost exclusively against
labels produced by SZZ, using protocols that ignore both commit order
and verification latency. The field therefore has \textbf{no measurement
of what SZZ-induced label noise costs a model under realistic deployment
conditions}, and no way to distinguish a model's ability to find defects
from its ability to reproduce the heuristic that labelled them.
\end{quote}

Four things must be separated to make progress: the noise in the labels;
the inflation contributed by the evaluation protocol; the
\emph{mechanism} by which noise damages an online learner; and whether a
learner can defend itself.

\textbf{This report covers the first two.} Phase 1 characterises the
noise; Phase 2 measures what it costs downstream and separates that cost
from the inflation contributed by the evaluation protocol. The remaining
two are stated as research questions, their designs are given in Chapter
6, and \textbf{no result is claimed for either.}

\section{1.5 Research gap}\label{research-gap}

Three lines of work approach this territory without closing it.

\textbf{SZZ evaluation research} measures how accurately SZZ variants
identify defect-introducing commits, including against
developer-informed oracles. It establishes that the labels are noisy. It
stops at the labels and does not measure the downstream cost.

\textbf{Online JIT-SDP research} addresses class imbalance evolution and
verification latency, and Cabral and Minku's ORB includes a safety
mechanism against potentially noisy minority-class examples. This work
treats noise as a nuisance to be dampened rather than as a measurable
quantity with a known, asymmetric structure and a known origin.

\textbf{Learning with noisy labels} provides loss correction, sample
selection and confidence-based methods. These are almost entirely batch
methods developed for symmetric or uniform label noise, and do not
transfer directly to a stream in which labels arrive late, arrive wrong,
and are later corrected.

\textbf{The gap is the intersection.} No existing work measures the cost
of \emph{SZZ-origin} noise --- asymmetric, variant-dependent, with known
flip rates --- under \emph{latency-aware online} evaluation, and no work
isolates which of SZZ's two error types an online learner actually
cannot survive. This report occupies that intersection, and its central
methodological device is the one that makes the intersection measurable:
\textbf{holding out an independently constructed reference label set so
that a model trained on SZZ can be scored against something other than
SZZ.}

\section{1.6 Research questions}\label{research-questions}

\textbf{RQ1.} How much do SZZ variants disagree with one another and
with an independently constructed reference, and what is the direction
of each variant's labelling bias? \emph{(Chapter 4)}

\textbf{RQ2.} How does the choice of SZZ label source shift measured
JIT-SDP performance across naive, chronological, and
prequential-with-latency evaluation regimes? \emph{(Chapter 5)}

\textbf{RQ3.} Through what mechanism does label noise degrade an online
learner --- specifically, how does ORB's oversampling response behave
under symmetric versus SZZ-calibrated asymmetric noise as the dose
increases? \emph{(Not addressed in this report; design in \S{}6.2)}

\textbf{RQ4.} Can an online learner recover performance lost to
SZZ-origin noise without access to clean labels, and without sacrificing
performance when the labels are clean? \emph{(Not addressed in this
report; design in \S{}6.2)}

\section{1.7 Contributions}\label{contributions}

\textbf{A class-conditional characterisation of SZZ noise.} Six variants
scored over an identical 27,319-commit denominator: a precision ceiling
of 27.2\%, bifurcated bias (\ensuremath{\rho}\ensuremath{_0} 0.067--0.263, \ensuremath{\rho}\ensuremath{_1} 0.359--0.733), and
inter-variant agreement of \ensuremath{\kappa} up to 0.933 against reference agreement of
\ensuremath{\kappa} 0.158--0.202. \textbf{Variants share failure modes, so their mutual
agreement is not evidence of correctness.} The flip-rate pairs are
released as a calibration artefact.

\textbf{A separation of protocol inflation from label effects.} Random
\emph{k}-fold inflates a high-capacity model by +0.137 MCC (21 of 21
projects); self-scoring inflates it by a further +0.230. The two
protocol effects together are roughly nine times the label-source
effect, and the inflation scales with model capacity.

\textbf{A calibrated noise model released as an artefact.} The
class-conditional flip rates (\ensuremath{\rho}\ensuremath{_0}, \ensuremath{\rho}\ensuremath{_1}) for all six variants, estimated
over a common denominator and serialised for reuse, so that downstream
noise studies can calibrate against measured SZZ behaviour rather than
assume symmetric noise.

\textbf{A reproducible pipeline with an integrity gate.} The full study
runs on a continuous-integration runner from two small committed files,
behind an assertion that recomputes the training data's confusion matrix
and refuses to spend compute if it disagrees with the reported noise
rates.

\section{1.8 Thesis organisation}\label{thesis-organisation}

\textbf{Chapter 2} surveys JIT-SDP, SZZ and its variants, evaluation
methodology, online learning under latency, and learning with noisy
labels, and positions this work against each.

\textbf{Chapter 3} states the research design: the phase pipeline, the
corpus and its reference labels, the SZZ toolchain, the metrics and
statistical protocol, and the validity framework --- stated in advance
rather than assembled afterwards.

\textbf{Chapters 4 and 5} report Phases 1 and 2, answering RQ1 and RQ2
respectively.

\textbf{Chapter 6} concludes: it synthesises the two results, gives the
full treatment of validity threats, and sets out the design of the
remaining phases together with what this report establishes that they
now depend on.

\textbf{A note on how the results chapters are written.} One finding in
this report contradicts a claim made earlier in the same project: a
decomposition of the verification-latency effect, which was withdrawn
after it was found not to be identified. The superseded claim is
\textbf{stated, together with the design error that produced it}, rather
than removed --- see \S{}5.3. The correction is part of the evidence: it
records which conclusions survived contact with a control and which did
not.
